\documentclass{article}
\usepackage[T1]{fontenc}
\usepackage{lmodern}
\usepackage{graphicx}
\usepackage{amsmath,amssymb}
\usepackage[a4paper,margin=2cm]{geometry}
\usepackage[absolute,overlay]{textpos}
\setlength{\TPHorizModule}{1mm}
\setlength{\TPVertModule}{1mm}
\usepackage[indonesian]{babel}
\usepackage{hyperref}
\setlength{\emergencystretch}{2em}
% unit: Halaman 23

\begin{document}

% === Halaman 23 (python/native) ===

Elliptic Curve Discrete Logarithm Problem 

(ECDLP)

• Menghitung kP = Q mudah, tetapi menghitung k jika diketahaui P dan Q adalah 

sulit. Inilah ECDLP yang menjadi dasar ECC.

• ECDLP dirumuskan sebagai berikut:

 Diberikan P dan Q adalah dua buah titik di kurva eliptik, carilah integer k 

sedemikian sehingga Q = kP

• Secara komputasi sulit menemukan k, jika k adalah bilangan yang besar.  Nilai k 

adalah logaritma diskrit dari Q dengan basis P.     *)

• Pada algoritma ECC, Q adalah kunci publik, k adalah kunci privat, dan P 

sembarang titik pada kurva eliptik.

23

Catatan: ingatlah kP = Pk , sehingga Q = kP = Pk, k adalah logaritma diskrit dari Q

\end{document}
